\documentclass[10pt,a4paper]{amsart}
\usepackage[margin=19mm]{geometry}
\usepackage{amsmath,amssymb,mathtools}
\usepackage[scr=rsfs,cal=euler]{mathalfa}
\usepackage{xfrac}
\usepackage[unicode,hidelinks,bookmarks=false]{hyperref}
\newcommand{\ie}{i.e.\ }
\newcommand{\cf}{cf.\ }
\newcommand{\resp}{respectively\ }
\newcommand{\Subsection}{\subsection}
\providecommand{\nobreakdash}{\nobreak\hbox{-}}
\setlength{\emergencystretch}{2em}
\multlinegap0pt
\usepackage{bm}
\usepackage{bbm}
\usepackage{dsfont}
\usepackage{xspace}
\usepackage{cancel}
\usepackage{mathtools}
\usepackage{amsthm}
\makeatletter
\@ifundefined{lemma}{\newtheorem{lemma}{Lemma}}{}
\@ifundefined{proposition}{\newtheorem{proposition}{Proposition}}{}
\makeatother
\newcommand{\AP}{\textnormal{\textbf{AP}}}
\newcommand{\R}{\mathbb{R}}
\newcommand{\boA}{\mathcal{A}}
\newcommand{\boI}{\mathcal{I}}
\newcommand{\boK}{\mathcal{K}}
\newcommand{\boS}{\mathcal{S}}
\newcommand{\boW}{\mathcal{W}}
\newcommand{\dq}{\dot{q}}
\title[Action-minimizing periodic orbits of the Lorentz force equat]{Action-minimizing periodic orbits of the Lorentz force equation with dominant vector potential}
\author{Manuel Garz\'on, Salvador L\'opez-Mart\'inez}
\date{Source: arXiv:2510.25462v1 (CC BY 4.0)}
\begin{document}
\maketitle

\noindent\emph{Source and licence.} Action-minimizing periodic orbits of the Lorentz force equation with dominant vector potential. Version arXiv:2510.25462v1 (CC BY 4.0), distributed under the Creative Commons Attribution 4.0 International licence, as recorded in the arXiv licence metadata for this exact version and shown on the abstract page https://arxiv.org/abs/2510.25462v1. Full rights notice: licenses/arxiv-2510.25462-CC-BY-4.0.txt.\par\smallskip

\noindent\textit{Editorial scope.} Counted unit: the Proposition whose source label is \texttt{MainTheorem2'} in the article (source0.tex lines 649--651, source label \texttt{MainTheorem2'}), delivered together with the article's own complete original proof of it (source0.tex lines 653--708, one \texttt{proof} environment of 56 lines). No mathematics is written, repaired or re-derived: statement and proof are the article's own text line for line. Required context retained uncounted: cond:varphi0 (lines 223--225); cond:varphi1 (lines 226--228); cond:varphi2 (lines 230--232); eq:poincare (lines 325--327); lemma:negative (lines 590--592); def:Jac (lines 599--601); def:g(t) (lines 604--606); def:p-eps (lines 612--614); ineq:epsilon1 (lines 632--634). Every definition, display and label of those lines is retained unchanged. Omitted from this package and not redistributed: 1--222, 229--229, 233--324, 328--589, 593--598, 602--603, 607--611, 615--631, 635--648, 652--652, 709--866. Nothing inside a retained excerpt is omitted or altered: every retained body is delivered exactly as the article's lines are written, so no omission receipt of any kind accompanies this package. Citations: the retained excerpts carry no citation command at all, so no bibliography entry of the article is transcribed and no reference is redistributed. Reference adaptation: none was needed and none was made; every reference inside the delivered unit resolves inside it, so no unresolved reference or citation is produced. Printed slips kept literal: no printed word, symbol, hypothesis or number of the counted statement or of its proof was altered. Layout and class: the article's own class is \texttt{article} with packages including geometry, amsmath, amsthm, amscd, amssymb, amsfonts, graphics, xcolor, cancel, comment, hyperref, mathtools; the wrapper is the lane's pinned \texttt{amsart} editorial wrapper at 10pt on A4, which restates the theorem-like environments the retained text needs and adds the article's own macro definitions that the retained text needs (\texttt{\textbackslash AP}, \texttt{\textbackslash R}, \texttt{\textbackslash boA}, \texttt{\textbackslash boI}, \texttt{\textbackslash boK}, \texttt{\textbackslash boS}, \texttt{\textbackslash boW}, \texttt{\textbackslash dq}), with extra wrapper packages (bm, bbm, dsfont, xspace, cancel, mathtools). The delivered numbering is the unit's own. Assets: no retained span contains a figure environment, a graphics-inclusion command, a PDF-page inclusion, a table or a TikZ picture; no image file is copied into this package, the package declares no asset dependency and no image file is redistributed with it. Licence: version arXiv:2510.25462v1 of this article is distributed under the Creative Commons Attribution 4.0 International licence, as recorded in the arXiv licence metadata for this exact version and shown on the abstract page https://arxiv.org/abs/2510.25462v1. A full rights notice is delivered at licenses/arxiv-2510.25462-CC-BY-4.0.txt. Characters: every retained excerpt is pure ASCII, so the delivered engine, XeLaTeX with its default Latin Modern text font, reports no missing character.
\begin{equation}\label{cond:varphi0}
    \partial_t A(\cdot,b_n)\not\equiv 0,\quad\hbox{for all }n\in\mathbb{N},
\end{equation}

\begin{equation}\label{cond:varphi1}
	    \lim\limits_{n\to\infty}\varphi(b_n)\|\tilde A(\cdot,b_n)\|^{-2}_2=\lim\limits_{n\to\infty}\varphi(b_n)\left(\int_0^T\left|\tilde A(t,b_n)\right|^2dt\right)^{-1}=0,
\end{equation}

\begin{equation}\label{cond:varphi2}
	\lim_{n\to\infty}\max\{|\nabla\Phi(t,y)|:\,t\in[0,T],\,|y|\geq |b_n|-T\}\|\tilde A(\cdot,b_n)\|^{-1}_2=0.
\end{equation}

\begin{equation}\label{eq:poincare}
	\frac{\pi^2}{T^2}\|\tilde q\|^2_2\leq\|\dq\|^2_2,\quad\text{for every }q\in\boW,
\end{equation}

\begin{lemma}\label{lemma:negative}
For any $A\in\boA$ there exists a point $p\in\boK$ such that $\boI_0(p)<0$.
\end{lemma}

\begin{equation}\label{def:Jac}
	M=\|\nabla A(t,x)\|_{L^\infty(\R^4)}.
\end{equation}

\begin{equation}\label{def:g(t)}
	g(t)=\int_0^t \tilde A(s,b) ds,\ \text{for every }t\in\R.
\end{equation}

\begin{equation}\label{def:p-eps}
  p=-\varepsilon \tilde g + b,\quad\hbox{for all }\varepsilon>0,
\end{equation}

    \begin{equation}\label{ineq:epsilon1}
        \boI_0(p)\leq-\varepsilon\|\dot{g}\|_2^2   \left(1-\varepsilon\frac{\pi+MT}{\pi}\right).
    \end{equation}

\begin{proposition}\label{MainTheorem2'}
    Let $(A,\Phi)\in\AP$ satisfy conditions \eqref{cond:varphi0}, \eqref{cond:varphi1} and \eqref{cond:varphi2}, for some sequence $\{b_n\}\subset\R^3$ with $\lim\limits_{n\to\infty}|b_n|=\infty$. Then there exists $q\in\boK_\Lambda$ such that $\boI(q)<0$.
\end{proposition}

\begin{proof}
    Let $\{b_n\}\subset\boS^c$ be the sequence with $\lim_{n\to\infty}|b_n|=\infty$ given by the statement. Similarly to \eqref{def:g(t)}, let us define the periodic function $g_n\in\boW$ as follows
\begin{equation}
    g_n(t)=\int_0^t\tilde A(s,b_n)ds,\quad \text{for all }n\in\mathbb{N}.
\end{equation}
Recall that $A(t,b_n)$ is non-constant for all $n\in\mathbb{N}$, and so
\begin{equation}
    \dot g_n(t)=\tilde A(t,b_n)\not\equiv 0,\quad\hbox{for all }n\in\mathbb{N}.
\end{equation}
Additionally, for $\varepsilon>0$, let $p_n$ be defined analogous to \eqref{def:p-eps}:
\begin{equation}
    p_n=-\varepsilon\tilde g_n + b_n,\quad\hbox{for all }n\in\mathbb{N}.
\end{equation} 
Since $\tilde A(t,b_n)\in\tilde\boW$, for all $n\in\mathbb{N}$, it follows that
\begin{equation}
    \max\limits_{t\in[0,T]}|\tilde A(t,b_n)|\leq T\|\partial_t A(t,x)\|_{L^\infty(\R^4)}\leq TM,\quad\hbox{for all }n\in\mathbb{N},
\end{equation}
where $M$ is given in \eqref{def:Jac}. Therefore, the sequence of norms $\|\dot g_n\|_\infty$ is uniformly bounded, and $\varepsilon>0$ can be chosen sufficiently small so that $p_n\in\boK$ for all $n\in\mathbb{N}$. Consequently, proceeding as in the proof of Lemma~\ref{lemma:negative}, we obtain the inequality
    \begin{equation}\label{ineq:epsilon2}
        \boI(p_n)\leq-\varepsilon\|\dot{g}_n\|_2^2   \left(1-\varepsilon\frac{\pi+MT}{\pi}\right)-\int_0^T\Phi(t,p_n)dt,\quad\hbox{for all }n\in\mathbb{N},
    \end{equation}
    which trivially reduces to \eqref{ineq:epsilon1} when $\Phi \equiv 0$. Hence, we now focus on estimating the term involving $\Phi$. First, notice that, for any $s\in [0,1]$, we have 
\begin{equation*}
    |s p_n(t)+(1-s)b_n|\geq |b_n|-\varepsilon| \tilde g_n|,\quad\hbox{for all }n\in\mathbb{N}. 
\end{equation*}
Then, recalling that 
\begin{equation}
    \|\tilde g_n\|_\infty\leq T\|\dot g_n\|_\infty,\quad\hbox{and}\quad\varepsilon \|\dot g_n\|_\infty<1,\quad\hbox{for all }n\in\mathbb{N},
\end{equation}it follows that
\begin{equation}\label{ineq:tvm}
    |s p_n(t)+(1-s)b_n|\geq |b_n|-T.
\end{equation}
    In particular, for all $n$ large enough, independent of $t$, every point in the segment joining $p_n(t)$ to $b_n$ belongs to $\boS^c$. Thus, by the mean value theorem there exists $s_n(t)\in [0,1]$ such that
\begin{align*}
    \left|\int_0^T(\Phi(t,p_n)-\Phi(t,b_n))dt\right|&\leq \varepsilon\int_0^T|\tilde g_n||\nabla\Phi(t,s_n(t)p_n+(1-s_n(t))b_n)|dt
    \\
    &\leq\frac{\varepsilon T}{\pi}\|\dot g_n\|_2\left(\int_0^T \left|\nabla\Phi(t,s_n(t)p_n+(1-s_n(t))b_n)\right|^2dt\right)^\frac12,
\end{align*}
where \eqref{eq:poincare} is used in the last step. As a consequence, recalling \eqref{ineq:tvm}, we deduce
\begin{equation*}
    \left|\int_0^T(\Phi(t,p_n)-\Phi(t,b_n))dt\right|\leq \frac{\varepsilon T^2}{\pi}\|\dot g_n\|_2\max\{|\nabla\Phi(t,x)|:\, t\in [0,T],\, |x|\geq|b_n|-T\}.
\end{equation*}
Therefore, condition \eqref{cond:varphi2} allows us to take $n$ sufficiently large so that
\begin{equation*}
    \left|\int_0^T(\Phi(t,p_n)-\Phi(t,b_n))dt\right|\leq \varepsilon^2\|\dot g_n\|_2^2.
\end{equation*}
Substituting this into \eqref{ineq:epsilon2} yields
\begin{equation}\label{ineq:epsilon2'}
        \boI(p_n) \leq-\varepsilon\|\dot{g}_n\|_2^2   \left(1-\varepsilon\frac{2\pi+MT}{\pi}\right)- \int_0^T\Phi(t,b_n)dt,
    \end{equation}
    for all $n\in\mathbb{N}$ sufficiently large. In addition, since $\int_0^T\Phi(t,b_n)dt=\varphi(b_n)$, condition \eqref{cond:varphi1} implies that
\begin{equation*}
        \boI(p_n) \leq-\varepsilon\|\dot{g}_n\|_2^2   \left(1-\varepsilon\frac{3\pi+MT}{\pi}\right),
    \end{equation*}
    for all $n \in \mathbb{N}$ sufficiently large. Therefore, taking $\varepsilon < \left(3 + MT\pi^{-1}\right)^{-1}$, we obtain that $\boI(p_n) < 0$, for all $n \in \mathbb{N}$ sufficiently large.
\end{proof}

\end{document}
